\section{PCS As A Concrete MLPCS Interface}

In this explainer, a multilinear PCS is the concrete implementation level for
the same evaluation-handle interface that the native MLPCS model idealizes.
The native object is a tensor
\[
  h\in\F^{\B^d}
\]
and its multilinear extension
\[
  \MLE(h):\F^d\to\F.
\]
The ideal MLPCS functionality stores \(h\) and answers evaluation requests
\[
  (\widehat O_h,r)\longmapsto \MLE(h)(r).
\]
A PCS replaces that ideal storage by a binding commitment and an opening proof:
after commit time, the verifier should be able to check claims of the form
\[
  \MLE(h)(r)=v
\]
without reading all \(2^d\) entries of \(h\).

In our notation, the nontrivial opening primitive is \(\EncMLE[C]\), not
MLPCS itself.  Given an encoded-MLE opening protocol for a code
\[
  C:\F^{\B^d}\to\F^{D},
\]
the ordinary MLPCS interface is obtained directly:
commit by fixing the encoded string \(y=C(h)\) behind a coordinate handle
\(O_y\), and open an evaluation claim \(\MLE(h)(r)=v\) by running
\(\EncMLE[C]\) on \((O_y,r,v)\).  Equivalently, an MLPCS handle is a
coordinate handle together with the declared code \(C\); each evaluation
opening is an encoded-MLE proof against that same fixed handle.

This is why we do not use a separate ideal \(\mathcal F_{\mathrm{PCS}}\) in the
rest of the paper.  When the text says ``PCS,'' it means a cryptographic
realization of the MLPCS evaluation-handle interface, usually through encoded
oracle strings, \(\EncMLE\) openings, proximity checks, coordinate openings,
Merkle authentication, and Fiat--Shamir.  When the text needs the ideal object,
it names
\(\mathcal F_{\mathrm{MLPCS}}\) or \(\mathcal F_{\mathrm{HMLPCS}}^\delta\)
directly.

\subsection{Representation Matters}

A concrete PCS must choose what tensor it commits to.  In this paper, a tensor
symbol such as \(y\) or \(h\) means whichever representation is being used at
that point, and we will say what it represents before using an evaluation
formula.

Two representations are especially common.  If \(y\in\F^{\B^d}\) stores Boolean
evaluations of a multilinear polynomial \(f:\F^d\to\F\), then
\[
  y[b]=f(b)
  \qquad\text{and}\qquad
  f=\MLE(y).
\]
In that case
\[
  f(z)
  =
  \sum_{b\in\B^d} y[b]\chi_b(z),
\]
where
\[
  \chi_b(z)
  =
  \prod_{i=0}^{d-1}\bigl(b_i z_i+(1-b_i)(1-z_i)\bigr).
\]
If instead \(y\) stores multilinear coefficients, then
\[
  f(X_0,\ldots,X_{d-1})
  =
  \sum_{b\in\B^d} y[b]\prod_{i=0}^{d-1}X_i^{b_i}.
\]
In that case
\[
  f(z)
  =
  \sum_{b\in\B^d} y[b]\prod_{i=0}^{d-1}z_i^{b_i}.
\]
These are different tensors related by a linear change of basis.  The protocol
goal \(f(z)=v\) is the same, but the public weights in the sum are different.
We should therefore never infer \(f=\MLE(y)\) unless \(y\) has just been declared to
be the Boolean-evaluation tensor.

\subsection{Trivial Realization}

The simplest concrete realization of the MLPCS evaluation-handle interface is
to reveal the entire representation tensor to the verifier.
Figure~\ref{fig:protocol-trivial-pcs} is useful as a baseline: it is correct,
but it gives up succinctness completely.

\begin{protocolbox}{Protocol: Trivial PCS}{fig:protocol-trivial-pcs}
\noindent\(\underline{\mathsf{TrivialPCS.Commit}(\prover:h\in\F^{\B^d},\ \verifier:\bot)\to C_h}\)

\smallskip
\begin{enumerate}[leftmargin=0.24in,label=\textbf{T\arabic*.},itemsep=0.08em,topsep=0.1em,parsep=0pt]
  \item \(\prover\to\verifier\): send the full tensor \(h\).
  \item \(\verifier\): store \(h\) under a fresh handle \(C_h\).
  \item \(\verifier\): return \(C_h\).
\end{enumerate}

\smallskip
\noindent\(\underline{\mathsf{TrivialPCS.Eval}(\prover:\bot,\ \verifier:(C_h,r,v))\to\{0,1\}}\)

\smallskip
\begin{enumerate}[leftmargin=0.24in,label=\textbf{U\arabic*.},itemsep=0.08em,topsep=0.1em,parsep=0pt]
  \item \(\verifier\): retrieve the stored tensor \(h\) under \(C_h\).
  \item \(\verifier\): compute \(u\gets \MLE(h)(r)\).
  \item \(\verifier\): if \(u\ne v\), return reject.
  \item \(\verifier\): return accept.
\end{enumerate}
\end{protocolbox}

The protocol in Figure~\ref{fig:protocol-trivial-pcs} is correct.  It is also
exactly what we are trying to avoid.  It leaks the whole polynomial
representation and makes the verifier store or scan \(2^d\) field elements.  For
the discussion here, privacy is not the central problem; efficiency and binding
are.  Later versions of polynomial commitments may also hide $f$, but the
Blaze/BaseFold construction is already interesting before adding hiding.
